% TOP_PROBLEM: {"id": 3150, "title": "Does the chromatic symmetric function distinguish between trees?", "category_id": 3, "category": {"id": 3, "name": "graph_theory", "display_name": "Graph Theory", "description": "Problems involving graphs, networks, and their properties.", "slug": "graph-theory", "order_index": 3, "created_at": "2026-07-31T15:26:25.671Z"}, "status": "open", "problem_number": "OPG-36880", "set_id": 12}
% FIRST_POSTED: 2026-10-01
% ENTRY_KIND: statement_only
% UnsolvedMath ID 3150; source catalogue code OPG-36880
% !TeX program = lualatex
% Definition and sources only; this file is not an LLM solution attempt.
\documentclass[11pt,a4paper]{article}
\usepackage[margin=25mm]{geometry}
\usepackage{fontspec}
\setmainfont{lmroman10-regular.otf}[BoldFont=lmroman10-bold.otf,
 ItalicFont=lmroman10-italic.otf,BoldItalicFont=lmroman10-bolditalic.otf]
\usepackage{amsmath,amssymb,mathtools}
\usepackage{unicode-math}
\setmathfont{latinmodern-math.otf}
\AtBeginDocument{\let\setminus\smallsetminus}
\usepackage{xurl,hyperref}
\hypersetup{colorlinks=true,urlcolor=blue}
\setcounter{secnumdepth}{0}
\setlength{\parindent}{0pt}
\setlength{\parskip}{5pt}
\setlength{\emergencystretch}{3em}
\begin{document}
\section*{Does the chromatic symmetric function distinguish between trees?}\label{3150}
{\small UnsolvedMath ID 3150 · OPG-36880 · Graph Theory · OpenGarden}
\subsection{Definitions and mathematical statement}
Let $G=(V,E)$ be a finite simple graph and let
$x_1,x_2,\ldots$ be commuting indeterminates. A proper colouring by
positive integers is a function $f:V\to\mathbb Z_{>0}$ with
$f(u)\neq f(v)$ whenever $uv\in E$. Define the chromatic symmetric
function by the formal sum
\[
 X_G(x_1,x_2,\ldots)
   =\sum_{f\text{ proper}}\prod_{v\in V}x_{f(v)}.
\]
This is a homogeneous symmetric function of degree $|V|$. Every
particular monomial has a finite coefficient, so no convergence
of the infinite sum is required. Its coefficient records the
number of proper colourings with the corresponding sizes of
colour classes.

A tree is a finite connected graph without a cycle. The question
asks whether two nonisomorphic trees can have the same chromatic
symmetric function. In its affirmative-distinguishability form,
the target assertion is
\[
 X_T=X_{T'}\quad\Longrightarrow\quad T\cong T'
 \qquad\text{for all finite trees }T,T'.
\]
Equality means coefficientwise equality in the same indeterminates;
isomorphism means a vertex bijection preserving adjacency. Equality
automatically forces the trees to have the same order, by
homogeneity. The question therefore concerns information beyond
the number of vertices. Specializing the first $q$ variables to
one and all others to zero counts proper $q$-colourings, but the
full symmetric function retains the separate colour-class sizes
discarded by that specialization.
\subsection{Short English statement}
Does the full chromatic symmetric function distinguish every pair of nonisomorphic finite trees?
\subsection{Sources}
\begin{itemize}
\item[S1] ULAM.AI and the UnsolvedMath Contributors, supplied problems.json, record 3150 (OPG-36880), OpenGarden collection. Catalogue identity and source transcription; CC BY 4.0.
\url{https://huggingface.co/datasets/ulamai/UnsolvedMath}.
\item[S2] Open Problem Garden, Does the chromatic symmetric function distinguish between trees?. Original problem and discussion; the mathematical formulation below is expanded from the supplied source record.
\url{http://www.openproblemgarden.org/op/does_the_symmetric_chromatic_function_distinguish_trees}.
\end{itemize}
\end{document}
